\SourceLeaf{259}%
dem linea brevissima \(DB\), pro quo, manente \(p\), characteres \(q',r',\varphi',\psi',S'\) eadem designent, quae \(q,r,\varphi,\psi,S\) pro puncto \(B\). Ita oritur triangulum inter puncta \(A,B,C\), cuius angulos per \(A,B,C\), latera opposita per \(a,b,c\), aream per \(\sigma\) denotamus; mensuram curvaturae in punctis \(A,B,C\) resp. per \(\alpha,\beta,\gamma\) exprimemus. Supponendo itaque (quod licet), quantitates \(p,q,q-q'\) esse positivas, habemus
\[
\begin{gathered}
A=\varphi-\varphi',\quad B=\psi,\quad C=\pi-\psi',\\
a=q-q',\quad b=r',\quad c=r,\quad\sigma=S-S'
\end{gathered}
\]

Ante omnia aream \(\sigma\) per seriem exprimemus. Mutando in [7] singulas quantitates ad \(B\) relatas in eas, quae ad \(C\) referuntur, prodit formula pro \(S'\), unde, usque ad quantitates sexti ordinis obtinemus
\[
\begin{aligned}
\sigma=\tfrac{1}{2}p(q-q')\{&1-\tfrac{1}{6}f^0(pp+qq+qq'+q'q')\\
&-\tfrac{1}{60}f'p(6pp+7qq+7qq'+7q'q')\\
&-\tfrac{1}{20}g^0(q+q')(3pp+4qq+4qq'+4q'q')\}
\end{aligned}
\]
Haec formula, adiumento seriei [2] puta
\[
c\sin B=p(1-\tfrac{1}{3}f^0qq-\tfrac{1}{4}f'pqq-\tfrac{1}{2}g^0q^3-\text{etc.})
\]
transit in sequentem
\[
\begin{aligned}
\sigma=\tfrac{1}{2}ac\sin B\{&1-\tfrac{1}{6}f^0(pp-qq+qq'+q'q')\\
&-\tfrac{1}{60}f'p(6pp-8qq+7qq'+7q'q')\\
&-\tfrac{1}{20}g^0(3ppq+3ppq'-6q^3+4qqq'+4qq'q'+4q'^3)\}
\end{aligned}
\]

Mensura curvaturae pro quovis superficiei puncto fit (per art. 19, ubi \(m,p,q\) erant quae hic sunt \(n,q,p\))
\[
\begin{aligned}
&=-\frac{1}{n}\mathord{.}\frac{\mathrm{d}\mathrm{d}n}{\mathrm{d}q^2}=-\frac{2f+6gq+12hqq+\text{etc.}}{1+fqq+\text{etc.}}\\
&=-2f-6gq-(12h-2ff)qq-\text{etc.}
\end{aligned}
\]
Hinc fit, quatenus \(p,q\) ad punctum \(B\) referuntur,
\[
\beta=-2f^0-2f'p-6g^0q-2f''pp-6g'pq-(12h^0-2f^0f^0)qq-\text{etc.}
\]
nec non
\[
\begin{aligned}
\gamma&=-2f^0-2f'p-6g^0q'-2f''pp-6g'pq'-(12h^0-2f^0f^0)q'q'-\text{etc.}\\
\alpha&=-2f^0
\end{aligned}
\]
